%!TeX root=MemoriaTFG.tex

\chapter{Conclusiones}
\label{cap:conclusiones}

\section{Respuesta a la pregunta de investigación}
\label{sec:respuesta}

La pregunta formulada en sección~\ref{sec:pregunta} se respondió mediante
la evaluación empírica de catorce modelos fundacionales sobre doce
datasets externos públicos armonizados a una ontología jerárquica
L1/L2/L3, bajo cuatro protocolos principales (sondeo lineal,
ajuste fino supervisado, segmentación binaria promptable y
clasificación zero-shot multimodal) y tres protocolos
complementarios (equidad por fototipo cutáneo, comparación con
modelos de lenguaje multimodales generalistas, auditoría de
solapamiento con el dataset de preentrenamiento Derm1M). La
evidencia disponible sostiene cuatro afirmaciones cualitativas, en
las que se resumen los hallazgos defendibles del trabajo.

\paragraph{Primera. Ventaja del preentrenamiento sobre el dominio
clínico.}
PanDerm Large constituye, entre los encoders contrastados con
pesos públicos, el modelo con mejor rendimiento agregado sobre los
diez datasets externos contrastados. La mejora media respecto a
PanDerm Base es de $+4{,}1$\,pp en accuracy, $+9{,}0$\,pp en
BAcc y $+3{,}6$\,pp en AUROC bajo sondeo lineal. La mejora
absoluta es máxima sobre Dermnet ($+10{,}0$\,pp accuracy),
WSI~patches ($+8{,}7$\,pp) y PAD-UFES-20 ($+4{,}7$\,pp),
datasets con mayor número de clases o mayor heterogeneidad de
adquisición. El ajuste fino con la receta del paper original
aumenta la BAcc en $+47{,}1$\,pp sobre HAM10000 y
$+19{,}4$\,pp sobre PAD-UFES-20, atribuible al aprendizaje de las
clases minoritarias. La especialización del dominio se cuantifica
también frente a arquitecturas convolucionales supervisadas: la
ventaja de PanDerm Large sobre la mejor CNN crece con la dificultad
del dataset, desde $+0{,}5$\,pp accuracy sobre HAM10000 hasta
$+8{,}2$\,pp sobre PAD-UFES-20.

\paragraph{Segunda. Saturación temprana del sondeo lineal.}
El sondeo lineal sobre encoders preentrenados específicos de
dominio alcanza saturación con un volumen de datos etiquetados muy
reducido: con $1\%$ del conjunto de entrenamiento de HAM10000
($82$ imágenes), PanDerm Large recupera $95{,}7\%$ de la
accuracy y $96{,}2\%$ del AUROC del modelo entrenado con el
conjunto completo; con $5\%$ ($410$ imágenes), AUROC alcanza
$97{,}7\%$ del valor máximo. Sobre PAD-UFES-20, $14$ imágenes
de entrenamiento ($1\%$) bastan para alcanzar $0{,}907$
AUROC, equivalente al $95{,}5\%$ del rendimiento con el conjunto
completo. Esta eficiencia es relevante para contextos clínicos con
disponibilidad limitada de etiquetas: enfermedades raras, centros
de bajo volumen asistencial, subgrupos poco representados y
extensión a especialidades dermatológicas no cubiertas por los
benchmarks habituales.

\paragraph{Tercera. Brecha entre modelos específicos y LLMs
multimodales generalistas.}
La comparación de paradigmas sobre HAM10000 cuantifica un
gradiente de rendimiento de $43{,}5$\,pp de accuracy entre el
mejor especialista ajustado (PanDerm Base FT con TTA, $0{,}920$) y el
mejor modelo de lenguaje multimodal generalista
\emph{zero-shot} (GPT-4o, $0{,}485$). La adaptación con LoRA
sobre MedGemma~27B reduce parcialmente esta brecha
($+13{,}7$\,pp Acc HAM, $+18{,}3$\,pp Acc DDI), pero no
sistemáticamente entre datasets ($-1{,}7$\,pp sobre
PAD-UFES-20). La consecuencia operativa para sistemas de apoyo al
diagnóstico es que los LLMs multimodales generalistas no
constituyen sustitutos defendibles de los encoders fundacionales
específicos del dominio en el régimen clínico evaluado, sin
perjuicio de su papel complementario como generadores de
razonamiento textual o explicaciones.

\paragraph{Cuarta. \emph{Gap} sistemático por fototipo cutáneo.}
Los cinco encoders contrastados sobre Fitzpatrick17k completo
presentan \emph{gap}~I-VI positivo sobre la formulación de tres
clases de malignidad, con valores en BAcc entre $+0{,}124$
(BiomedCLIP) y $+0{,}306$ (DINOv2 ViT-L/14). La paradoja
BiomedCLIP, con \emph{gap} mínimo asociado a la peor BAcc global
($0{,}590$), ilustra que la métrica de equidad aislada no
constituye un criterio decisional: un modelo uniformemente inferior
puede presentar baja desigualdad sin ser equitativo en sentido
sustantivo. El modelo con mejor compromiso entre rendimiento global
y equidad es PanDerm Large (BAcc $0{,}699$, \emph{gap}
$+0{,}217$); DermLIP v2 ofrece la mejor BAcc global
($0{,}721$) con \emph{gap} superior ($+0{,}281$).

\paragraph{Síntesis. No hay modelo universalmente superior, sino
especialización por tarea.}
Leídas en conjunto, las cuatro afirmaciones anteriores no convergen
en la designación de un modelo vencedor único, sino en una lectura
matizada de la pregunta de investigación. La evidencia empírica
apunta a que el rendimiento dermatológico no es una propiedad
absoluta del codificador, sino una función conjunta del nivel
ontológico evaluado, del régimen de datos etiquetados disponibles,
de la modalidad de entrada y de la estrategia de adaptación. Cada
escenario favorece una combinación distinta de modelo y protocolo:
clasificación general, subcategoría clínica L2, cola larga L3,
ranking probabilístico, segmentación y cribado de seguridad de
melanoma no se resuelven con el mismo método óptimo. Esta lectura,
desarrollada como tesis transversal en
sección~\ref{sec:disc-especializacion}, desplaza el cuello de
botella del avance: deja de residir únicamente en la capacidad
representacional del encoder y pasa a localizarse en la taxonomía
diagnóstica, la calidad y cobertura de las etiquetas, el
alineamiento texto-imagen, la formulación de los \emph{prompts}, la
equidad por fototipo cutáneo y la estrategia de despliegue. La
respuesta a la pregunta formulada en sección~\ref{sec:pregunta} es,
por tanto, condicional: el modelo y el protocolo adecuados se
seleccionan en función del objetivo clínico y de las restricciones
de cómputo, datos y supervisión, no en abstracto.

\paragraph{Corolario: la palanca de mejora son los datos.}
Si el cuello de botella ya no reside únicamente en la capacidad
representacional del codificador, el corolario transversal del trabajo
es que la palanca dominante de mejora son los datos. La evidencia lo
sugiere de forma recurrente: la ventaja de PanDerm y de DermLIP
procede de la escala y el dominio de su preentrenamiento; la
adaptación eficiente de Dermapixel R0 sólo es viable porque dispone de
datos clínicos en castellano; el nivel L3 fracasa por falta de
muestras en la cola larga; y la línea de mayor impacto potencial es la
ampliación del dataset con el banco hospitalario del HUSLL. El
siguiente salto en rendimiento clínico no se espera, por tanto, de
sustituir el codificador, sino de ampliar y mejorar los datos ---en
cobertura diagnóstica, representación de fototipos y registro
lingüístico--- sobre los que se adapta y evalúa.

\section{Contribuciones reproducibles}
\label{sec:contribuciones}

Las contribuciones del trabajo se enumeran en orden decreciente
de centralidad respecto a la pregunta de investigación.

\subsection{Contribuciones en el cuerpo del documento}
\label{sec:contrib-cuerpo}

\begin{enumerate}
\item \textbf{Reproducción experimental de PanDerm sobre
arquitectura aarch64.} El pipeline de PanDerm Base y Large se ha
verificado sobre NVIDIA DGX Spark con chip Grace Hopper GB10,
$128$\,GB de memoria unificada, CUDA~13.0 y precisión BF16, con
reproducción de los rangos numéricos reportados por
Yan et~al.~\cite{panderm2025} sobre los datasets compartidos. La
portabilidad al chip Grace Hopper queda documentada como base
operativa para grupos de investigación con hardware no estándar.

\item \textbf{Benchmark homogéneo de catorce modelos fundacionales
sobre doce datasets.} El estudio comparativo aplica cuatro
protocolos principales bajo configuración estándar
(sección~\ref{sec:protocolos}) sobre PanDerm Base y Large, DINOv2,
tres variantes de DermLIP, BiomedCLIP, SigLIP-Large, SAM2.1-Large,
GPT-4o, Gemini~2.5~Pro, MedGemma~4B Vision y~27B Text, y BLIP-2, más
dos líneas de base de redes convolucionales supervisadas
(ConvNeXt-Large y EfficientNetV2-Large) como referencia no
fundacional. Las cifras se reportan
con desglose por dataset y métrica
(tabla~\ref{tab:lp-panderm}, tabla~\ref{tab:benchmark},
tabla~\ref{tab:llm-comparativa}).

\item \textbf{Ontología jerárquica L1/L2/L3.} Cuatro categorías
etiológicas, $43$ subcategorías clínicas y $367$
diagnósticos individuales codificados en SNOMED~CT y CIE-10,
validados por revisión experta de la Dra.~R.~Taberner. La
ontología se aplica a la armonización de once de los doce
datasets externos contrastados, agregando $72\,654$ imágenes
sobre vocabulario común y habilitando comparaciones inter-dataset
defendibles a nivel L1 y L2.

\item \textbf{Dataset DermapixelAI~1.0.} Primer dataset
dermatológico verificado en español con imagen, metadatos y texto
narrativo asociado ($1\,089$ imágenes, $669$ casos clínicos,
mediana de $223$ palabras por caso). Particionado bajo regla
\emph{case-aware} ($891 / 157 / 41$) y distribuido bajo
licencia CC-BY-NC-SA~4.0 con datasheet formal
(Anexo~\ref{cap:anexoc}) y documento de entrega
(Anexo~\ref{cap:anexoe}).

\item \textbf{Auditoría sistemática de solapamiento con Derm1M.}
Procedimiento por hash MD5 exhaustivo sobre los doce datasets
externos contrastados, con identificación de Dermnet como único
conjunto con solapamiento exacto ($100\%$ contenido en Derm1M) y
declaración de \emph{caveat} explícito sobre las cifras zero-shot de
la familia DermLIP en ese dataset.

\item \textbf{Marco metodológico reproducible de evaluación.} Más allá
de las contribuciones anteriores tomadas por separado, su articulación
bajo un protocolo común ---benchmark homogéneo, ontología jerárquica,
auditoría de \emph{leakage}, análisis de equidad por fototipo,
interpretabilidad mediante Sparse Autoencoders, segmentación
promptable y clasificación zero-shot multimodal--- constituye en sí
misma un marco de evaluación reproducible y transferible a futuros
modelos fundacionales dermatológicos, con código, \emph{splits}
canónicos y \emph{outputs} documentados en el
Anexo~\ref{cap:anexoa}.
\end{enumerate}

\subsection{Contribuciones complementarias}
\label{sec:contrib-anexos}

\begin{enumerate}
\setcounter{enumi}{5}
\item \textbf{Análisis empírico de equidad por fototipo.}
Evaluación de cinco encoders sobre Fitzpatrick17k completo
($16\,577$ imágenes, $6$ fototipos cutáneos) con
cuantificación del \emph{gap}~I-VI por modelo
(tabla~\ref{tab:equidad-fototipo}) y análisis cruzado
\emph{Light}~$\leftrightarrow$~\emph{Dark}
(tabla~\ref{tab:cross-eval}).

\item \textbf{Caracterización de la sensibilidad al tokenizador en
DermLIP.} Identificación del tokenizador y la formulación del
\emph{prompt} como factores críticos del rendimiento zero-shot,
con diferencia de $48{,}8$\,pp AUROC entre la peor y la mejor
configuración sobre HAM10000.

\item \textbf{Adaptación de SAM2.1 al dominio dermatoscópico
mediante \emph{fine-tuning} del decodificador.} Ajuste denso del
\emph{mask\_decoder} y el \emph{promptable encoder} ($\sim 4{,}2$\,M
pesos entrenables), congelando el
\emph{image\_encoder} Hiera-Large preentrenado. Resultado de Dice
$0{,}947$ sobre ISIC2018 ($+2{,}6$\,pp sobre la cifra
publicada por Yan et~al.~\cite{panderm2025}), con transferencia
fuera de dominio prácticamente nula sobre ISIC2017 ($0{,}945$)
y PH2 ($0{,}960$).

\item \textbf{Ensemble de seguridad para detección de melanoma.}
Combinación OR \emph{top-3} de PanDerm Large ajustado sobre
HAM10000 (M1) y SigLIP-Large SO400M con sondeo lineal, que alcanza
\emph{recall} de melanoma del $100\%$
($70$ de $70$, \emph{top-3}) sobre el conjunto de test de
HAM10000.

\item \textbf{Modelo de cuello de botella conceptual SAE${}\times{}$SkinCon.}
Replicación del mecanismo de interpretabilidad por Sparse
Autoencoders sobre PanDerm Large (SAE Large,
$1\,024 \to 16\,384$ \emph{features}, entrenado sobre
$50\,454$ imágenes) y su validación frente al diccionario
clínico SkinCon~\cite{skincon} mediante un \emph{Concept
Bottleneck Model}, con AUROC media de $0{,}880$ sobre los
$22$ conceptos con cobertura suficiente y $11$ conceptos con
AUROC del mejor predictor superior a $0{,}80$. Los experimentos
E1--E4 de validación conceptual sobre Derm7pt se detallan en el
Anexo~\ref{cap:anexof}.

\item \textbf{Dermapixel R0: primer modelo fundacional
dermatológico en español.} Modelo inicializado desde el
codificador PanDerm Large y adaptado sobre el dataset
DermapixelAI~1.0, articulado en una rama supervisada (cabeza L2
con LoRA $r=16$, BAcc $0{,}363 \pm 0{,}007$ y AUROC
$0{,}855 \pm 0{,}006$, sección~\ref{sec:dermapixel-spanderm-v0})
y una rama contrastiva multimodal castellano
(sección~\ref{sec:spanderm-clip}). Constituye una versión base
entrenada sobre un corpus reducido: la contribución reside en ser
el primero de su clase y en la metodología reproducible que
establece, no en un rendimiento competitivo. Su escalado con el
banco hospitalario del HUSLL se documenta como línea de
continuación en la sección~\ref{sec:linea-spanderm}.
\end{enumerate}

\section{Hallazgos metodológicos transferibles}
\label{sec:hallazgos-conc}

Los hallazgos metodológicos del trabajo son transferibles a
escenarios clínicos análogos y constituyen aportaciones operativas
más allá del benchmark concreto evaluado. El primero es de orden
estructural ---la inexistencia de un método óptimo único--- y
enmarca a los restantes, que precisan las condiciones bajo las que
cada estrategia resulta preferible.

\paragraph{Ausencia de un método óptimo único.}
El hallazgo transversal del trabajo es que ningún codificador ni
estrategia de adaptación domina todos los escenarios: la
configuración preferible depende del nivel ontológico, del régimen
de datos, de la modalidad de entrada y del objetivo clínico
(clasificación general, subcategoría clínica L2, cola larga L3,
ranking probabilístico, segmentación o cribado de seguridad). Esta
especialización por tarea, sostenida en
sección~\ref{sec:disc-especializacion}, implica que el diseño de un
sistema de apoyo dermatológico debe formularse como selección
condicional de modelo y protocolo, y no como adopción de un único
encoder de referencia.

\paragraph{Régimen de aplicabilidad del sondeo lineal.}
La comparación cuantitativa sobre cuatro datasets
(sección~\ref{sec:disc-lp-vs-ft}) delimita un régimen operativo: para
datasets con $N_{\mathrm{train}} \lesssim 500$ imágenes, el
sondeo lineal sobre PanDerm Large supera al ajuste fino completo
en accuracy, BAcc, AUROC y F1 ponderada. El umbral se identifica
sobre DDI ($N_{\mathrm{train}} = 427$), donde el sondeo lineal
mantiene AUROC $0{,}827$ frente a $0{,}682$ del ajuste fino.

\paragraph{Compatibilidad tokenizador-encoder en zero-shot.}
La AUROC en clasificación zero-shot con modelos vision-lenguaje
contrastivos depende críticamente de la compatibilidad entre el
tokenizador empleado para los \emph{prompts} y el espacio de
\emph{embeddings} del encoder textual del modelo. La verificación
empírica de esta compatibilidad sobre un conjunto de validación
del dominio es prerrequisito de cualquier despliegue zero-shot
defendible.

\paragraph{Ajuste del decodificador sobre encoders promptables
preentrenados a gran escala.}
La combinación de SAM2.1-Large preentrenado con \emph{fine-tuning}
del decodificador de máscaras, manteniendo congelado el encoder
visual, supera al
\emph{fine-tuning} completo de baselines dermatológicos
específicos, con actualización de aproximadamente $4{,}2$\,M
parámetros entrenables (inferior al $2\%$ del total). La
generalización fuera de dominio es prácticamente nula, lo que
sugiere que las primitivas geométricas codificadas por el encoder
visual transfieren entre fuentes de adquisición sin
reentrenamiento.

\paragraph{Limitación operativa del \emph{gap} aislado como métrica
de equidad.}
La paradoja BiomedCLIP ---\emph{gap}~I-VI mínimo ($+0{,}124$ en
BAcc) asociado a la peor BAcc global ($0{,}590$)--- demuestra que la
métrica \emph{gap} aislada puede inducir conclusiones erróneas
sobre la equidad sustantiva de un modelo. El reporte conjunto de
BAcc global, AUROC global y \emph{gap} es metodológicamente
preferible y debería adoptarse como práctica estándar en estudios
de equidad dermatológica.

\paragraph{Sondeo lineal supera al ajuste fino bajo sobreajuste.}
El caso DDI ($N_{\mathrm{train}} = 427$, descenso de AUROC
$-14{,}5$\,pp con ajuste fino respecto al sondeo lineal) es
diagnóstico de la posibilidad de sobreajuste del encoder ajustado
pese a las augmentaciones de la receta canónica (Mixup, CutMix,
\emph{weighted random sampler}). La heurística operativa derivada
es que el ajuste fino sobre datasets pequeños debe ir acompañado
de monitorización de AUROC sobre conjunto de validación
independiente, y descartarse si el ranking de probabilidades
degrada.

\paragraph{Recuperación $k$-NN como alternativa competitiva en la
cola larga.}
La evidencia empírica apunta a que, sobre el nivel diagnóstico fino
L3, dominado por una distribución de cola larga estricta, la
recuperación por vecinos más próximos con índice FAISS sobre
\emph{embeddings} congelados resulta competitiva frente a las
cabezas paramétricas, que carecen de muestras suficientes por clase
para estimar fronteras de decisión estables
(sección~\ref{sec:disc-especializacion}). La heurística operativa
es que, cuando la cardinalidad de clases excede la disponibilidad de
etiquetas, la recuperación no paramétrica sobre el espacio de
representación es preferible a la clasificación supervisada directa.

\paragraph{Interpretabilidad conceptual emergente vía Sparse
Autoencoders.}
La descomposición del espacio de representación de PanDerm Large
mediante Sparse Autoencoders revela \emph{features} alineadas con
conceptos dermatológicos clínicamente interpretables sin supervisión
conceptual durante el preentrenamiento, validadas frente al
diccionario SkinCon~\cite{skincon} y frente a las anotaciones
\emph{Seven-Point Checklist} de Derm7pt. Esta interpretabilidad
emergente habilita modelos de cuello de botella conceptual que
acotan el compromiso interpretabilidad-precisión, y constituye una
vía de trazabilidad clínica transferible a otros encoders
fundacionales del dominio.

\section{Aportaciones a la práctica clínica}
\label{sec:aportaciones-clinicas}

El trabajo aporta cuatro elementos directamente aplicables al
diseño de sistemas de apoyo al diagnóstico dermatológico en el
contexto del Servicio de Dermatología del HUSLL y, por extensión,
al circuito de teledermatología institucional descrito en
sección~\ref{sec:trabajo-previo}.

\paragraph{Arquitectura técnica defendible.}
Encoder PanDerm Large preentrenado y congelado por defecto, con
cabeza lineal multi-clase mediante regresión logística L-BFGS
($C = 1{,}0$, \texttt{max\_iter}~$=5\,000$) en datasets con
$N_{\mathrm{train}} \lesssim 500$, o cabeza ajustada con la
receta de PanDerm Large en datasets con
$N_{\mathrm{train}} \gtrsim 5\,000$ y desbalance acusado.
Segmentación con SAM2.1-Large adaptado mediante \emph{fine-tuning}
del decodificador cuando la tarea requiera localización de lesión. Ensemble OR \emph{top-3}
de M1 (PanDerm Large \emph{fine-tuned}) y SigLIP-Large con sondeo
lineal cuando la tarea sea cribado oportunista de melanoma con
prioridad de \emph{recall}.

\paragraph{Ontología armonizada como columna vertebral.}
La ontología L1/L2/L3 con codificación SNOMED~CT y CIE-10
proporciona una columna vertebral semántica reutilizable para
sistemas de información hospitalarios y compatible con la
nomenclatura clínica del entorno. La presentación jerárquica al
usuario clínico (predicción L3 acompañada de su agregación L2 y
L1) aporta robustez frente a la incertidumbre en la cola larga sin
sacrificar especificidad cuando ésta sea fiable.

\paragraph{Dataset DermapixelAI~1.0 como recurso público en español.}
Habilita la reutilización del archivo Dermapixel en proyectos
académicos posteriores, la adaptación de modelos vision-lenguaje a
texto clínico en castellano y la validación de sistemas de apoyo
sobre material de referencia docente reconocida en el ámbito
hispanohablante.

\paragraph{Caveats explícitos sobre fuga de información.}
La auditoría de solapamiento entre datasets de evaluación y
preentrenamiento aporta una práctica de reporte transferible a
otros benchmarks dermatológicos. La declaración explícita de
\emph{caveats} de leakage en las tablas
(tabla~\ref{tab:datasets-resumen}) y la consistencia entre la cifra
reportada y su \emph{caveat} asociado debería adoptarse como
condición de defensibilidad en publicaciones del campo.

\section{Líneas de investigación abiertas}
\label{sec:lineas-abiertas}

Las líneas abiertas se organizan en cuatro bloques. Conviene distinguir
dos registros: parte del trabajo aquí recogido ya está \emph{ejecutado}
y reportado en el capítulo~\ref{cap:resultados} ---la cobertura
experimental sobre DermapixelAI~1.0, el eje Derm7pt + Sparse Autoencoders
y la selección de arquitectura de segmentación---, y se presenta como
base consolidada; el resto ---escalado con el banco HUSLL, validación
clínica prospectiva, aprendizaje activo, IA agéntica y reentrenamiento
multilingüe de la rama lingüística--- es agenda futura genuina,
condicionada a recursos o aprobaciones específicas. Todas descansan sobre
las contribuciones reproducibles de la sección~\ref{sec:contribuciones}.

\subsection{Escalado del dataset con el banco hospitalario HUSLL}
\label{sec:linea-husll}

La extensión de DermapixelAI~1.0 con el banco hospitalario del HUSLL
($\sim 22\,000$ estudios en archivo, con imagen clínica y dermatoscópica
cuando disponible, diagnóstico final, edad y sexo) es la línea de mayor
impacto potencial sobre el rendimiento clínico real. Sus prerrequisitos
son cuatro: aprobación del Comité Ético de Investigación Clínica de las
Illes Balears (CEIC-IB), un protocolo de anonimización conforme al RGPD y
la Ley Orgánica~$3/2018$, un esquema de licencia compatible con la
naturaleza hospitalaria del material (probablemente más restrictivo que
CC-BY-NC-SA~4.0), y un equipo de validación experta liderado por la
Dra.~R.~Taberner con apoyo de residentes del HUSLL. Su incorporación
habilita la cabeza L3 directa de Dermapixel R0
(sección~\ref{sec:linea-spanderm}), la validación clínica prospectiva
(sección~\ref{sec:linea-prospectivo}) y un análisis de equidad por
fototipo más representativo de la población atendida por el sistema
sanitario español que el disponible en Fitzpatrick17k.

\paragraph{Aprendizaje activo \emph{in-domain} sobre la cola larga L3.}
\label{sec:linea-active-learning}
La trilogía contrastiva de Dermapixel R0
(sección~\ref{sec:spanderm-clip}) mostró que ampliar el entrenamiento con
escala genérica externa (v5) no rellena la cola larga L3 ---$327$ clases
con $\leq 5$ imágenes--- y erosiona la alineación cross-modal castellana,
mientras que la iteración v4 ya alcanza el umbral \emph{a priori} de
\emph{accuracy}@1 $\sim 0{,}25$. La palanca correcta no es más datos
genéricos, sino aprendizaje activo dirigido \emph{in-domain}: identificar
las clases L3 peor cubiertas e invertir en ellas anotación experta, con
candidatos generados por el prototipo sobre cohortes del HUSLL bajo
revisión de la Dra.~R.~Taberner, hasta superar las $\sim 30$ imágenes por
clase que activan la cabeza L3 directa
(sección~\ref{sec:linea-spanderm}). La ablación del Anexo~\ref{cap:anexoi}
lo refuerza: la palanca pendiente es anotación \emph{in-domain}, no escala
genérica.

\paragraph{Auditoría sistemática de \emph{leakage} extensible.}
\label{sec:linea-leakage}
La auditoría por hash MD5 desarrollada para Derm1M
(sección~\ref{sec:auditoria}) es directamente extensible a los demás
datasets de preentrenamiento del ecosistema (SkinGPT-4, MONET, MM-Skin,
SkinFM, entre otros): aplicar el mismo procedimiento sobre los
\emph{splits} de evaluación canónicos y publicar las cardinalidades de
intersección, con su script auditable, permitiría recalibrar las cifras
\emph{zero-shot} publicadas sobre esos datasets.

\subsection{Dermapixel R0: adaptación dermatológica multimodal en
español}
\label{sec:linea-spanderm}

\textbf{Dermapixel R0} constituye, a juicio de este trabajo, su
aportación más original. Hasta donde alcanza la revisión bibliográfica
realizada en este trabajo, es una de las primeras adaptaciones
dermatológicas multimodales con texto clínico narrativo en español
construidas sobre un modelo fundacional abierto. La denominación se emplea
en sentido preciso: no es un preentrenamiento fundacional a gran escala
---como PanDerm o Derm1M---, sino la adaptación de un modelo fundacional
preexistente (PanDerm Large) al registro clínico en castellano; el
calificativo \emph{primero} se acota a esa categoría ---pesos abiertos,
texto clínico narrativo y registro en español--- y no a una prioridad
histórica absoluta. Su relevancia no reside en una métrica de máximo
rendimiento ---es, por diseño, una versión base (R0) sobre un dataset
reducido, un punto de partida y no un modelo terminado---, sino en lo que
demuestra: que un modelo dermatológico especializado en castellano puede
construirse sobre un modelo fundacional liberado mediante adaptación
eficiente, con un presupuesto de cómputo y datos al alcance de una
institución hospitalaria y un grupo académico.

Dermapixel R0 se inicializa desde PanDerm Large, se adapta sobre
DermapixelAI~1.0 y se articula en dos ramas complementarias. La rama
\textbf{supervisada} ejecuta una cabeza L2 ($38$ clases) con LoRA $r=16$
($524\,$K parámetros, $0{,}17\%$ del encoder), que reporta
BAcc~$=0{,}363 \pm 0{,}007$ y AUROC~$=0{,}855 \pm 0{,}006$ sobre L2
($+19{,}5$\,pp BAcc sobre el sondeo lineal congelado, AUROC en empate
estadístico con DermLIP v2), con ranking L3 por prototipos coseno
(Acc@1~$=0{,}167$; sección~\ref{sec:dermapixel-spanderm-v0}). La cabeza L3
\emph{directa} ($251$ diagnósticos) queda condicionada al banco HUSLL: su
activación exige $\geq 30$ imágenes por clase, umbral que DermapixelAI~1.0
no satisface en aislamiento ($178$ entradas L3 con $\leq 5$ imágenes). La
rama \textbf{zero-shot multimodal} (dual-encoder estilo CLIP, encoder
visual PanDerm Large + textual \texttt{multilingual-e5-base}) ejecuta
cinco iteraciones (sección~\ref{sec:spanderm-clip}): v4 alcanza L3
\emph{accuracy}@1 $=0{,}2475$ (umbral \emph{a priori} $0{,}25$), v5
documenta una ablación negativa de escala ($0{,}2079$) y v2 logra
recuperación cross-modal con \emph{lift} de $\mathbf{14{,}0\times}$
($i2t$ R@$5=0{,}654$), con calibración por \emph{temperature
scaling}~\cite{guo2017} que reduce el ECE L3 un $97\%$. Ambas ramas operan
en regímenes de uso distintos (vocabulario abierto frente a cerrado) y se
integran como módulos paralelos del prototipo (Anexo~\ref{cap:anexoh}). En
conjunto, Dermapixel R0 traslada al dominio una conclusión más amplia: la
frontera de la especialización ya no la marca el acceso a grandes
infraestructuras de preentrenamiento, sino la disponibilidad de datos
clínicos bien curados y de un protocolo de adaptación adecuado, replicable
por otras instituciones con archivos clínicos en su propia lengua.

\paragraph{Cobertura experimental sobre DermapixelAI~1.0: técnicas
ejecutadas y pendientes.}
\label{sec:dermapixel-futuro}
Las secciones~\ref{sec:dermapixel-results} y~\ref{sec:dermapixel-zs}
evalúan los codificadores sobre DermapixelAI~1.0 como extractores
congelados con cabezas ligeras o clasificación \emph{zero-shot}, siguiendo
el protocolo estándar de la literatura que cuantifica limpiamente el valor
del preentrenamiento. La tabla~\ref{tab:dermapixel-pendiente} resume el
estado de las técnicas de \emph{transfer learning} sobre el dataset: el
ajuste fino denso, la adaptación LoRA de Dermapixel R0, el \emph{Test-Time
Augmentation}, el ensemble multi-codificador y la pérdida focal están
\emph{ejecutados} y reportados en el capítulo~\ref{cap:resultados};
permanecen pendientes la línea CBM sobre conceptos SkinCon (condicionada a
anotación, Anexo~\ref{cap:anexof}) y la extensión vía banco HUSLL
(bloqueada por la revisión CEIC). Bajo aplicación combinada de las
técnicas pendientes y la ampliación del dataset hasta superar el umbral
L3, la cifra esperable para un clasificador desplegable sobre L2 en
castellano se estima en BAcc entre $0{,}45$ y $0{,}55$; es una estimación
prospectiva, no un compromiso, y queda condicionada a la disponibilidad
efectiva del banco hospitalario.

\begin{table}[H]
\centering
\footnotesize
\setlength{\tabcolsep}{4pt}
\renewcommand{\arraystretch}{0.95}
\caption{Técnicas de \emph{transfer learning} pendientes de
aplicación sobre DermapixelAI~1.0 con el cuerpo experimental del
presente trabajo. \emph{Coste}: orden de magnitud sobre NVIDIA DGX
Spark GB10. \emph{Esperado}: mejora cualitativa estimada a partir
de las referencias internas del trabajo. \emph{Estado}: situación
respecto al cierre del presente TFG.}
\label{tab:dermapixel-pendiente}
\begin{tabular}{p{0.19\textwidth}p{0.09\textwidth}p{0.15\textwidth}p{0.30\textwidth}p{0.15\textwidth}}
\hline
Método & Tipo & Coste estimado & Mejora esperada & Estado \\
\hline
Ajuste fino completo del encoder        & Supervisado            & $\sim 5$ min Spark (3 niveles) & $+18{,}2$\,pp BAcc L1 sobre LP CV; $-4{,}0$\,pp BAcc L2 frente a Dermapixel R0 LoRA; sin mejora sobre LP en L3 & \textbf{Ejecutado} (L1$+$L2$+$L3), sección~\ref{sec:dermapixel-ft} \\
LoRA $r=16$ (Dermapixel R0)             & PEFT                   & $\sim 12$ min Spark (3 seeds) & $+19{,}5$\,pp BAcc L2 sobre LP CV; AUROC L2 $0{,}855$ empata con DermLIP individual & \textbf{Ejecutado}, sección~\ref{sec:dermapixel-spanderm-v0} \\
\emph{Test-time augmentation}           & Inferencia             & $5$--$10$ min              & $+3$ a $+22$\,pp Acc@3 (cf.\ sección~\ref{sec:ensemble-results}) & \textbf{Probado}, sección~\ref{sec:dermapixel-tta} \\
Ensemble Base$+$Large$+$DermLIP         & Inferencia             & Trivial                    & $+6{,}7$\,pp AUROC L2 con DermLIP individual; $+2{,}1$\,pp AUROC L1 con avg Large$+$DermLIP & \textbf{Probado}, sección~\ref{sec:dermapixel-ensemble} \\
Pérdida focal / re-ponderación de clases & Entrenamiento         & $\sim 5$ min Spark         & $+3{,}3$\,pp BAcc L1 y $+5{,}9$\,pp BAcc L2 sobre LP CE ($\gamma = 2{,}0$); sin efecto L3 & \textbf{Ejecutado} ($\gamma \in \{1,\,2,\,5\} \times \{\text{none},\,\text{balanced}\}$), sección~\ref{sec:dermapixel-focal} \\
CBM sobre conceptos SkinCon              & Interpretable          & Anotación dermatológica    & Trazabilidad clínica                                          & Diseñado (Anexo~\ref{cap:anexof}) \\
Banco hospitalario HUSLL ($\sim 22$\,K) & Datos                  & $3$--$6$ meses (CEIC)      & Activación L3 directa (cabeza L3 de Dermapixel R0)            & Bloqueado (revisión ética) \\
\hline
\end{tabular}
\end{table}


\paragraph{Derm7pt + Sparse Autoencoders como ground-truth conceptual.}
\label{sec:linea-derm7pt-sae}
Las anotaciones \emph{Seven-Point Checklist} de Derm7pt (siete criterios
por imagen) son ground-truth conceptual no explotada por la literatura
previa sobre Sparse Autoencoders. Los cuatro experimentos de este eje
(correlación SAE-Derm7pt, \emph{Concept Bottleneck Model}, cruce con el
Grupo~A y \emph{fine-tuning} multitarea) están \emph{ejecutados}
(sección~\ref{sec:sae-derm7pt}, Anexo~\ref{cap:anexof}): la base
\emph{sparse} aprendida sobre Derm1M alinea \emph{features} con la
taxonomía dermatoscópica sin entrenamiento específico (AUROC efectivo
hasta $0{,}926$) y el \emph{fine-tuning} multitarea produce inferencia
simultánea de melanoma y los siete criterios en una pasada (AUROC
$0{,}891$). La línea genuinamente abierta es extender el diccionario
conceptual a los $15$ conceptos adicionales propuestos por la
Dra.~R.~Taberner (Anexo~\ref{cap:anexof}), condicionada a la anotación de
$\sim 200$ imágenes por concepto.

\subsection{Validación clínica prospectiva}
\label{sec:linea-prospectivo}

La validación clínica prospectiva del prototipo DermApIxel
(Anexo~\ref{cap:anexoh}) sobre cohorte hospitalaria real del HUSLL
requiere un protocolo aprobado por CEIC-IB con cuatro elementos:
\emph{endpoints} primarios (\emph{recall} de melanoma, derivaciones
evitadas) y secundarios (concordancia inter-observador, tiempo de
diagnóstico); cálculo muestral basado en los \emph{recall} preliminares de
$\sim 0{,}97$ obtenidos sobre HAM10000 con M1+SigLIP; revisores
independientes para arbitraje; y un plan de análisis estadístico
preregistrado. La cifra de \emph{recall} del $100\%$ sobre melanoma
reportada en sección~\ref{sec:ensemble-results} es prueba de concepto
sobre material de archivo, no resultado clínico validado: su extrapolación
al flujo asistencial real exige esta validación, requisito de cualquier
despliegue operativo.

\paragraph{Recuperación visual densa como sistema de apoyo.}
\label{sec:linea-rag}
La recuperación visual densa basada en FAISS sobre los \emph{embeddings}
de Derm1M ($\sim 421\,000$ imágenes preindexadas) ya opera en el prototipo
(sección~\ref{sec:dermapixel-m4} del Anexo~\ref{cap:anexoh}) con latencia
inferior a $150$\,ms. La línea abierta es evaluar su valor clínico
mediante un estudio mixto sobre cohorte HUSLL (\emph{precision@k} sobre
L1/L2 y valoración de la utilidad de los casos recuperados), con la
integración en la historia clínica electrónica como prerrequisito técnico.

\paragraph{Selección de arquitectura para segmentación promptable.}
\label{sec:linea-seg-arquitectura}
La comparativa pareada de arquitecturas promptables
(sección~\ref{sec:seg-arquitecturas}) ya está ejecutada y delimita una
pauta de selección reproducible: la generación SAM2 domina sobre SAM~v1,
el factor determinante es la caja-\emph{prompt} y no el segmentador, y
MedSAM2-tiny (Dice~$0{,}9556$, $26$\,ms) es candidata de referencia para
despliegues con restricciones de cómputo. La línea abierta se reduce a dos
elementos: un detector de lesión extremo a extremo que sustituya la caja
manual y la validación prospectiva sobre adquisición clínica no
dermatoscópica.

\subsection{IA agéntica y modelos multimodales en español}
\label{sec:linea-agentica}

A medio plazo, la combinación de los componentes desarrollados (encoder
PanDerm, ontología L1/L2/L3, dataset DermapixelAI, recuperación visual
densa y LLM multimodal en local) habilita un agente clínico orquestado
bajo control del facultativo, que incorpore evidencia estructurada externa
y auditable. Frente al uso directo de un LLM multimodal generalista
(sección~\ref{sec:disc-llm}), este esquema reduciría el sesgo léxico
documentado y facilitaría la justificación clínica de la salida; su
desarrollo es condicional a pesos abiertos con capacidad de razonamiento
estructurado y se detalla en el Anexo~\ref{cap:anexoh}.

\paragraph{Reentrenamiento multilingüe de la rama lingüística.}
\label{sec:linea-multilingual}
La sensibilidad documentada al tokenizador (sección~\ref{sec:disc-zs})
sugiere que reentrenar el encoder textual de DermLIP sobre texto clínico
en español aportaría mejoras apreciables sobre el rendimiento
\emph{zero-shot} hispanohablante. La estrategia mínima viable aplica
\emph{continued pretraining} del encoder textual sobre los $669$ casos
clínicos de DermapixelAI~1.0 con su razonamiento narrativo, manteniendo
congelado el encoder visual; la ampliada requiere un dataset de pares
imagen-texto en español de mayor escala, condicional al escalado con el
banco HUSLL (sección~\ref{sec:linea-husll}).

\section{Cierre}
\label{sec:cierre}

El trabajo ha caracterizado empíricamente el rendimiento y la
capacidad de generalización de catorce modelos fundacionales para
imagen dermatológica disponibles, bajo cuatro protocolos
principales y tres complementarios, sobre doce datasets externos
armonizados a una ontología jerárquica de tres niveles validada
por revisión experta. La evidencia obtenida sostiene cuatro
afirmaciones cualitativas defendibles (sección~\ref{sec:respuesta}),
las contribuciones reproducibles enumeradas en
sección~\ref{sec:contribuciones}, los hallazgos metodológicos
transferibles de la sección~\ref{sec:hallazgos-conc} y cuatro
aportaciones a la práctica clínica
(sección~\ref{sec:aportaciones-clinicas}).

La lectura de conjunto de esta evidencia no es la designación de un
codificador vencedor, sino una tesis matizada: el rendimiento
dermatológico depende del nivel ontológico, del régimen de datos,
de la modalidad de entrada y de la estrategia de adaptación, de modo
que la elección de modelo y protocolo se resuelve en función del
objetivo clínico y no en abstracto
(sección~\ref{sec:disc-especializacion}). Esta lectura matiza
cualquier ordenación global de los modelos evaluados y desplaza la
palanca de mejora desde la capacidad representacional del encoder
hacia la taxonomía diagnóstica, la cobertura de etiquetas, el
alineamiento texto-imagen, los \emph{prompts}, la equidad por
fototipo cutáneo y la estrategia de despliegue.

Las líneas de investigación abiertas
(sección~\ref{sec:lineas-abiertas}) constituyen propuestas concretas con
experimentos identificados y prerrequisitos declarados, no trabajo
simplemente pendiente. Su ejecución completaría la cadena entre
caracterización empírica del estado del arte, construcción de
recursos públicos en español ---el benchmark DermapixelAI~1.0 y la
semilla Dermapixel R0, primer modelo fundacional dermatológico
adaptado al español---, validación clínica prospectiva y despliegue
operativo en el sistema sanitario, recorrido cuya primera etapa cierra
el presente documento.

La aportación intelectual del trabajo se sitúa en la conjunción de
tres elementos. El primero es una contribución empírica: la
evaluación rigurosa y multi-eje del estado del arte, que cuantifica
la dependencia del rendimiento respecto al nivel ontológico, el
régimen de datos, la modalidad y la estrategia de adaptación. El
segundo es una contribución de recurso: el dataset DermapixelAI~1.0
y la semilla Dermapixel R0
(contribuciones~\ref{sec:contrib-cuerpo}
y~\ref{sec:contrib-anexos}), primeros de su clase para texto
clínico dermatológico en español. El tercero es la formulación
explícita de líneas de continuación con base técnica y clínica
verificable. De esta evidencia se desprende que el avance del campo
no pasa por un único encoder de referencia, sino por una
arquitectura modular que combine codificadores, niveles ontológicos
y estrategias de adaptación según la tarea, junto con decisiones de
despliegue que contemplen cobertura de etiquetas, equidad y
trazabilidad. El siguiente paso natural es la primera publicación
derivada del trabajo, centrada en la comparativa de modelos
fundacionales bajo la ontología armonizada, seguida del escalado de
Dermapixel R0 como modelo fundacional dermatológico para texto
clínico en español a partir del dataset DermapixelAI~1.0 extendido
con el banco hospitalario del HUSLL.
